\DocumentMetadata{
  pdfversion=2.0,
  pdfstandard=ua-2,
  lang=en-US,
  tagging=on,
  tagging-setup={math/setup=mathml-SE,tabsorder=structure}
}
\documentclass[11pt,letterpaper]{article}
\usepackage[margin=0.68in,includefoot]{geometry}
\usepackage{newcomputermodern}
\usepackage{amsmath,amscd,mathtools}
\usepackage{tikz,adjustbox}
\providecommand{\square}{\mdlgwhtsquare}
\usepackage[hidelinks]{hyperref}
\hypersetup{
  pdftitle={Topology PhD Exam, May 2020},
  pdfauthor={University of Florida Department of Mathematics},
  pdfsubject={Graduate mathematics examination},
  pdfdisplaydoctitle=true,
  bookmarksopen=true
}
\setcounter{secnumdepth}{0}
\setlength{\parindent}{0pt}
\setlength{\parskip}{0.28em}
\setlength{\emergencystretch}{3em}
\setlength{\leftmargini}{2.15em}
\setlength{\leftmarginii}{2.65em}
\setlength{\leftmarginiii}{3em}
\renewcommand{\labelenumi}{\textbf{\theenumi.}}
\pagestyle{empty}
\raggedbottom
\allowdisplaybreaks
\newenvironment{examproblems}[1]{%
  \begin{enumerate}%
  \setcounter{enumi}{#1}%
  \setlength{\topsep}{0.35em}%
  \setlength{\partopsep}{0pt}%
  \setlength{\itemsep}{0.48em}%
  \setlength{\parsep}{0.18em}%
}{\end{enumerate}}
\newenvironment{examsubparts}{%
  \begin{enumerate}%
  \setlength{\topsep}{0.18em}%
  \setlength{\partopsep}{0pt}%
  \setlength{\itemsep}{0.16em}%
  \setlength{\parsep}{0.1em}%
}{\end{enumerate}}
\newenvironment{examinstructions}{%
  \par\begingroup\itshape
}{%
  \par\endgroup\vspace{0.18em}
}
\makeatletter
\renewcommand\section{\@startsection{section}{1}{0pt}%
  {-1.25ex plus -.25ex minus -.1ex}%
  {0.55ex plus .1ex}%
  {\normalfont\large\bfseries}}
\renewcommand{\@maketitle}{%
  \newpage\null\vskip -1em%
  {\footnotesize\bfseries
    \href{https://www.ufl.edu/}{University of Florida}, \href{https://math.ufl.edu/}{Department of Mathematics}\par}%
  \vskip -0.5em%
  \tagstructbegin{tag=Title}%
    {\LARGE\bfseries\href{https://gma.math.ufl.edu/exams/topology-phd/}{Topology PhD Exam}, May 2020\par}%
  \tagstructend%
  \vskip 0.25em%
  \tagmcbegin{artifact}%
    \hrule height 1pt%
  \tagmcend%
  \vskip 1em%
}
\makeatother
\title{Topology PhD Exam, May 2020}
\author{}
\date{}
\begin{document}
\maketitle
\thispagestyle{empty}
\begin{examinstructions}
Work the following problems and show all work. Support all statements to the best of your ability.
\end{examinstructions}
\begin{examproblems}{0}
\item
Let~\(X\) be a nonempty set. A nonempty collection~\(\Phi\) of subsets \(U\subseteq X\times X\) is a uniformity if it satisfies the following axioms:

\par
(1) If \(U\in\Phi\), then \(\Delta\subseteq U\), where~\(\Delta\) is the diagonal \(\{(x,x):x\in X\}\).

\par
(2) If \(U\in\Phi\) and \(U\subseteq V\subseteq X\times X\), then \(V\in\Phi\).

\par
(3) If \(U,V\in\Phi\), then \(U\cap V\in\Phi\).

\par
(4) If \(U\in\Phi\), then there is a \(V\in\Phi\) with \(V\circ V\subseteq U\), where \(A\circ B=\{(x,z)\mid \exists y\in X\text{ such that }(x,y)\in A\text{ and }(y,z)\in B\}\).

\par
(5) If \(U\in\Phi\), then \(U^{-1}\in\Phi\), where \(U^{-1}=\{(y,x)\mid (x,y)\in U\}\).

\par
The elements of~\(\Phi\) are called entourages. If \(x\in X\) and \(U\in\Phi\), set \(U[x]=\{y\mid (x,y)\in U\}\). A uniformity induces a topology on~\(X\): a set~\(O\) is open if for every \(x\in O\) there is an entourage~\(V\) such that \(V[x]\subseteq O\). A topological space is uniformizable if there is a uniformity compatible with the topology (that is, the topology generated by the uniformity is the topology on~\(X\)).

\par
Prove that a metric space~\(X\) is uniformizable.
\item
Let~\(T\) be the 2-dimensional torus. For each of the following, give an example or prove such an object does not exist.
\begin{examsubparts}
\item[\textnormal{(1)}]
A map \(f:T\to S^2\) of degree 1.
\item[\textnormal{(2)}]
A map \(g:S^2\to T\) of degree 1.
\end{examsubparts}
\item
The torus~\(T\) bounds a compact region~\(R\). Two copies of~\(R\) are glued together by identifying their boundaries via the map \(f:T\to T\) given by \(f(\theta,\varphi)=(\theta+\varphi,\theta+2\varphi)\) (here \(T=S^1\times S^1\), where points on \(S^1\) are identified with an angle \(\theta\in[0,2\pi)\)). The result is a compact 3-manifold~\(M\).
\begin{examsubparts}
\item[\textnormal{(1)}]
Compute the map \(f_*:H_1(T)\to H_1(T)\).
\item[\textnormal{(2)}]
Use the Mayer–Vietoris sequence to compute the homology of~\(M\).
\end{examsubparts}
\item
Let \(X\subset\mathbb{R}^3\) be the union of~\(n\) lines through the origin. Compute \(\pi_1(\mathbb{R}^3-X)\).
\item
Compute \(H^\bullet(M;\mathbb{Z})\), where~\(M\) is the space with homology groups 
\[
H_k(M;\mathbb{Z})=
\begin{cases}
\mathbb{Z}, & k=0,2,4,\\
\mathbb{Z}_3, & k=1,\\
\mathbb{Z}_2, & k=3,\\
0, & \text{otherwise}.
\end{cases}
\]
\end{examproblems}
\begin{examinstructions}
Answer the following with complete definitions, statements, or short proofs.
\end{examinstructions}
\begin{examproblems}{5}
\item
Prove that for a finite CW-complex~\(X\), \(H^1(X;\mathbb{Z})\) is free.
\item
Compute \(\chi(\mathbb{CP}^2\times\mathbb{RP}^4\times S^2)\).
\item
Give an example of a space that is connected but not path-connected.
\item
State the Urysohn Metrization Theorem.
\item
Prove that if \(m\ne n\), then \(\mathbb{R}^m\) is not homeomorphic to \(\mathbb{R}^n\).
\item
Let~\(X\) be a space and let \(X^*\) be its one-point compactification. Prove that if~\(X\) is locally compact and Hausdorff, then \(X^*\) is Hausdorff.
\item
State the Lefschetz Fixed Point Theorem.
\item
Describe all the connected covering spaces \(E\to S^1\times S^1\).
\item
Does the following exact sequence of abelian groups necessarily split? Prove or give a counterexample. 
\[
0\to A\to B\to\mathbb{Z}\to 0
\]
\item
Compute the integral homology of the space \(\mathbb{RP}^2\times\mathbb{CP}^3\).
\end{examproblems}
\end{document}
